\documentclass[11pt]{ctexart}
\usepackage[a4paper,margin=1in]{geometry}
\usepackage{graphicx}
\usepackage{xcolor}
\usepackage{hyperref}
\definecolor{refgray}{gray}{0.45}
\title{\textbf{市场最新观点汇总}\\[0.4em]{\large AI重塑就业与科技格局，私募股权进入结构性调整期}}
\author{KC桌面 自动生成}
\date{260705}
\begin{document}
\maketitle
\tableofcontents
\newpage
\section{一页摘要}
\begin{itemize}
  \item BCG指出未来2-3年美国50\%-55\%的岗位将被AI重塑，但大规模失业并非必然，软件工程师等岗位反而因需求扩张而增长。
  \item 麦肯锡2025技术趋势报告强调AI是激活其他技术的引擎，Agentic AI股权投资和岗位增长爆发。
  \item 麦肯锡全球私募市场报告揭示2025年交易额创新高但结构分化，行业从杠杆驱动转向运营创造价值。
\end{itemize}
\section{投行/券商}
今日暂无新增报告。

\begin{itemize}
  \item 今日暂无新增报告。
\end{itemize}
\section{战略咨询}
BCG与麦肯锡三份报告分别聚焦AI对就业的重塑、技术趋势的协同效应以及私募市场的结构性变化

\subsection{AI重塑就业：替代、增强与需求扩张的三重力量}
\textbf{AI对就业的影响并非简单的替代，而是通过替代、增强和需求扩张三种机制重塑岗位，软件工程师等岗位反而因AI降低生产成本、扩大需求而增加。}

\begin{itemize}
  \item BCG模型显示未来2-3年美国50\%-55\%的岗位将被AI重塑，但5年内真正被替代的岗位仅10\%-15\%，大部分员工保留岗位但工作方式彻底改变。
  \item 57\%的美国岗位自动化潜力低于40\%的阈值，这些岗位依赖物理在场、动手操作或人际互动，AI颠覆性影响有限。
  \item 需求弹性是关键变量：软件需求因AI降低成本而扩大，软件工程师岗位反而增加；而客服等结构化岗位需求有限，面临替代风险。
  \item 报告建议企业优先做增强而非替代，投资技能提升，避免因急于裁员而损失人才。
\end{itemize}
\begin{figure}[htbp]
\centering
\includegraphics[width=0.88\linewidth]{figures/fig_001.jpg}
\caption{Exhibit 1 - Over the next two to three years, 50\% to 55\% of jobs in the US will be reshaped by AI. For many employees, this will mean that they retain the same or a similar role but face radically new expectations for how they work and what they produce. For company leade}
\end{figure}
\begin{figure}[htbp]
\centering
\includegraphics[width=0.88\linewidth]{figures/fig_002.jpg}
\caption{Exhibit 2 - \# Task Automation Doesn’t Have to Mean Job Loss \#\# EXHIBIT 2 Agentic AI May Drive High Levels of Task Automation in 43\% of Jobs Sources: Revelio Labs; O\textbackslash{}*NET; US Bureau of Labor Statistics; BCG Henderson Institute analysis. Substitution Versus Augmentation. To}
\end{figure}
\begin{figure}[htbp]
\centering
\includegraphics[width=0.88\linewidth]{figures/fig_003.jpg}
\caption{Exhibit 3 - This dynamic is not new. Economists have long observed that efficiency improvements can increase total consumption rather than reduce it, a phenomenon often referred to as Jevons Paradox. When the cost of a resource falls, usage can rise. The same logic applie}
\end{figure}
{\small\color{refgray}\textbf{References:} [R001] 波士顿咨询：波士顿咨询揭示AI就业悖论，软件工程师反而因AI需求增长}

\subsection{2025技术趋势：AI作为激活器，Agentic AI爆发式增长}
\textbf{AI不再是孤立赛道，而是激活机器人、生物工程、能源等技术的引擎，Agentic AI从概念走向商业场景，股权投资和岗位增长显著。}

\begin{itemize}
  \item 麦肯锡追踪13项前沿技术，2024年股权投资在10个趋势中实现反弹，Agentic AI股权投资达11亿美元，岗位增长985\%。
  \item AI加速机器人训练、推动生物工程发现、优化能源系统，与其他技术组合产生最大价值，自主系统从试点走向实用。
  \item 规模化面临数据中心电力、供应链延迟、人才短缺等现实挑战，大国科技竞赛加剧。
  \item 报告首次将Agentic AI单独列为趋势，强调其从执行任务到学习、适应、协作的进化，人机关系从替代转向增强。
\end{itemize}
\begin{figure}[htbp]
\centering
\includegraphics[width=0.88\linewidth]{figures/fig_048.jpg}
\caption{Exhibit 1 - \#\# Wind farm maintenance crew Energy and sustainability Wind farms need to maintain peak efficiency year-round, maximizing clean-energy output while minimizing downtime and operational costs Immersive reality Technicians use augmented reality goggles that prov}
\end{figure}
\begin{figure}[htbp]
\centering
\includegraphics[width=0.88\linewidth]{figures/fig_049.jpg}
\caption{Exhibit 1 - Immersive reality Technicians use augmented reality goggles that provide visual guidance to help them maintain and repair complex turbine systems safely Advanced connectivity Low-Earth-orbit satellites provide technicians with access to real-time data and clou}
\end{figure}
\begin{figure}[htbp]
\centering
\includegraphics[width=0.88\linewidth]{figures/fig_050.jpg}
\caption{Exhibit 2 - McKinsey \& Company Note: Data includes private-market and public-market capital raises across venture capital and corporate and strategic M\&A (including joint ventures), private equity investments (including buyouts and private investment in public equity), an}
\end{figure}
{\small\color{refgray}\textbf{References:} [R003] 麦肯锡：麦肯锡2025技术趋势，AI不是主角，而是激活其他技术的引擎}

\subsection{私募市场结构性转变：从杠杆驱动到运营创造价值}
\textbf{2025年全球PE交易额破万亿美元创新高，但交易数量下降、持有期拉长、LP流动性枯竭成为结构性特征，行业进入需要真正运营能力的阶段。}

\begin{itemize}
  \item 2025年大额交易（超5亿美元）增长44\%至1.1万亿美元，但中小型buyout数量下降5\%，北美最明显。
  \item 超过16000家公司被持有超4年（历史最高），平均持有期超6.5年，LP的DPI占AUM仅6\%（2015-19年平均16\%）。
  \item 过去杠杆和估值扩张贡献59\%回报，未来需靠运营创造价值（operational alpha），AI正从试点走向生产，70\%的GP认为3-5年内影响翻天覆地。
  \item 超50亿美元大基金份额扩大，小而美基金缩水，专注特定赛道的specialist表现优于generalist，半流动性产品募资翻倍至2040亿美元。
\end{itemize}
\begin{figure}[htbp]
\centering
\includegraphics[width=0.88\linewidth]{figures/fig_008.jpg}
\caption{Exhibit 1 - A look at larger deals, a critical metric in assessing overall industry health, is revealing. For buyout and growth deals larger than \textbackslash{}\$500 million, 2025 was a record year, increasing 44 percent versus 2024 to \textbackslash{}\$1.1 trillion, higher than 2021's previous peak. }
\end{figure}
\begin{figure}[htbp]
\centering
\includegraphics[width=0.88\linewidth]{figures/fig_009.jpg}
\caption{Exhibit 2 - Geographical perspective sharpens the picture. In North America, buyout deal value increased 29 percent, accounting for 57 percent of total global buyout deal value (Exhibit 2). Europe also saw an increase (8 percent), but APAC declined (3 percent). This type }
\end{figure}
\begin{figure}[htbp]
\centering
\includegraphics[width=0.88\linewidth]{figures/fig_010.jpg}
\caption{Exhibit 3 - 1. Bigger deals are back: Deployment pressure is combining with buyers paying more for quality. The rebound in deal value occurred because acquirers are paying more, rather than doing more deals. The median EBITDA multiples that buyout firms paid reached a rec}
\end{figure}
{\small\color{refgray}\textbf{References:} [R002] 麦肯锡：私募股权已经驶入更硬的地面，不是所有车都能开过去}

\section{智库/国际机构}
今日暂无新增报告。

\begin{itemize}
  \item 今日暂无新增报告。
\end{itemize}
\section{结语}
今日国际信源主线清晰：AI正从概念走向系统性重塑就业与技术生态，而私募市场在交易额新高背后迎来结构性挑战，运营能力和技术整合成为未来竞争力的核心。
\newpage
\section{报告覆盖清单}
\begin{itemize}
  \item [R001] 战略咨询 | 战略咨询 / AI重塑就业：替代、增强与需求扩张的三重力量 - 波士顿咨询：波士顿咨询揭示AI就业悖论，软件工程师反而因AI需求增长
  \item [R002] 战略咨询 | 战略咨询 / 私募市场结构性转变：从杠杆驱动到运营创造价值 - 麦肯锡：私募股权已经驶入更硬的地面，不是所有车都能开过去
  \item [R003] 战略咨询 | 战略咨询 / 2025技术趋势：AI作为激活器，Agentic AI爆发式增长 - 麦肯锡：麦肯锡2025技术趋势，AI不是主角，而是激活其他技术的引擎
\end{itemize}
\section{Disclaimer}
{\scriptsize\color{refgray}
This community is a paid learning and information-sharing community created for general educational purposes only. The membership fee is charged solely for access to community discussions, educational content, organization, and operational costs. It is not an advisory fee, management fee, performance fee, brokerage fee, or compensation for personalized financial, investment, legal, tax, or professional advice.\par
I participate and share information solely as an individual and for educational discussion among community members. I am not acting as a financial adviser, investment adviser, broker, dealer, portfolio manager, legal adviser, tax adviser, accountant, fiduciary, or any other licensed professional adviser. Nothing shared in this community constitutes, and should not be interpreted as, financial advice, investment advice, legal advice, tax advice, a recommendation, solicitation, offer, endorsement, instruction, or guarantee to buy, sell, hold, short, trade, or invest in any security, cryptocurrency, fund, derivative, strategy, or other financial product.\par
All content shared in this community, including messages, discussions, links, files, charts, examples, market commentary, watchlists, AI-generated or AI-polished summaries, and third-party materials, is general, impersonal, and educational in nature. It does not take into account any individual's financial situation, investment objectives, risk tolerance, investment horizon, liquidity needs, tax status, legal restrictions, or suitability requirements. No content should be relied upon as suitable for any specific person, account, portfolio, or financial decision.\par
Information shared in this community may be incomplete, inaccurate, outdated, speculative, AI-generated, AI-edited, or based on third-party internet sources that have not been independently verified. I make no representation or warranty as to the accuracy, completeness, timeliness, reliability, or suitability of any information shared.\par
Financial markets involve substantial risk, including the possible loss of principal. Past performance, historical data, hypothetical examples, simulated results, backtests, personal opinions, or educational examples do not guarantee future results. Each member is solely responsible for conducting their own research, exercising independent judgment, and making their own decisions. Before making any financial, investment, legal, or tax decision, members should consult qualified and licensed professionals.\par
Participation in this community, payment of any membership fee, or communication with me or other members does not create any adviser-client, fiduciary, professional, contractual, agency, partnership, employment, or other advisory relationship. By joining or participating in this community, you acknowledge and agree that you are solely responsible for your own decisions, actions, transactions, profits, losses, risks, and consequences, and that I shall not be liable for any direct, indirect, incidental, consequential, financial, legal, tax, or other loss or damage arising from or related to any content, discussion, information, or material shared in this community.\par
}
\end{document}
